\section{Implementation checks}
\label{app:checks}
The pure-Python Adam was compared with \texttt{torch.optim.Adam} (PyTorch 2.14.0, CPU, float64,
\texttt{foreach}, \texttt{fused}, \texttt{capturable} and AMSGrad off, weight decay 0) on a 400-step
gradient sequence. The maximum relative differences were \ParityW{} (weights), \ParityM{} ($m$),
\ParityV{} ($v$) and \ParityUpd{} (updates), and the step counters were identical. Trace comparisons
further show that the following hold within numerical tolerance:
\begin{itemize}
\item \texttt{reset\_both} with decoupled counters equals a newly constructed torch optimizer;
\item \texttt{reset\_t} equals zeroing torch's \texttt{step};
\item the shared-keep variants equal zeroing \texttt{exp\_avg}/\texttt{exp\_avg\_sq} alone.
\end{itemize}
Decoupled single-moment resets have no stock-PyTorch equivalent and are checked only against their
explicit formulas. These are implementation checks, not proofs.

\section{Counterexample to scalar counter equivalence with \texorpdfstring{$\eps>0$}{eps>0}}
\label{app:counterexample}
Two coordinates receive the constant gradients $1$ and $10^{-8}$ for $10{,}000$ steps
($\beta_1{=}0.9$, $\beta_2{=}0.999$). From the same post-gradient $(m,v)$, one further step is taken with
the counter kept ($t{=}10{,}001$) and with the counter reset ($r{=}1$). With $\eps=0$ the ratio of the
reset update to the kept update is \CeEpsZeroRatioA{} for the first coordinate and \CeEpsZeroRatioB{}
for the second. This equals $\kappa(1,10{,}001)=\CeKappa$, which differs from the $t\to\infty$ value
\CeKappaInf{} by the finite-age factor. With $\eps=10^{-8}$ the ratios are \CeRatioA{} and \CeRatioB{}.
The two coordinates are therefore rescaled differently, and no single learning-rate multiplier reproduces
the counter reset. The implementation matches \eqref{eq:ratio} to a relative error of \CeMaxErr{}.
The same behaviour is observed with stock \texttt{torch.optim.Adam} when \texttt{step} is zeroed
(\texttt{tests/test\_torch\_parity.py}). This is a sanity check of elementary algebra, not a
contribution. Script: \texttt{analysis/counter\_equivalence.py}; output:
\texttt{runs/p4\_analysis\_counter/result.json}.

\section{Raw-trace precision of the \texttt{reset\_t} control}
\label{app:trace}
Only the first step after the intervention shares $(m,v)$ across branches, so one-step
ratios are checked at that step only, per tensor, over both conditions and three seeds.
\begin{itemize}
\item \texttt{reset\_t}/\keep: the maximum relative deviation from $\kappa(1,t{+}1)$ is \TraceResetTDev.
\item \texttt{reset\_both} shared-keep/decoupled: the ratio is predicted to be \TraceBothSkPred{} at
$\eps=0$; the maximum relative deviation is \TraceBothSkDev.
\item \texttt{reset\_v} shared-keep/decoupled: the ratio is predicted to be \TraceVSkPred{} at $\eps=0$,
but the observed ratios range from \TraceVSkMin{} to \TraceVSkMax.
\end{itemize}
For \texttt{reset\_v}, \TraceVSkEpsLimited{} of \TraceVSkTensors{} tensor--seed pairs have a ratio within
1\% of 1. Of the \TraceVSkBigCount{} pairs whose decoupled update norm exceeds 1, \TraceVSkBigNearOne{}
are within 1\% of 1 and none exceeds \TraceVSkBigMaxRatio, far below the $\eps=0$ prediction. By \eqref{eq:ratio}, a ratio of
1 is what results when $\eps$ dominates $\sqrt{\hat v}$ in both branches. The output bias, whose gradient
is never structurally zero, has the predicted ratio (at least \TraceVSkBiasMin). This is per-tensor
evidence consistent with near-zero-gradient coordinates driving the \texttt{reset\_v} blow-up. It is not
a coordinate-level verification (\textsc{checkpoint\_trace\_unavailable}).

Over the whole 200-step trajectories, \texttt{reset\_t} and its update-norm control differ by at most
\TrajLossDiff{} in training loss, \TrajProbeLossDiff{} in probe loss and \TrajThetaDiff{} in relative
parameter change, and by \TrajProbeAccDiff{} in probe accuracy. The ``$+0.000$'' in
Table~\ref{tab:interventions} is thus an exact tie in accuracy between trajectories that are not
bitwise identical. The small continuous differences are consistent with $\eps>0$ and rounding. The
control computes its own update direction and rescales only its per-tensor norm
(\texttt{osrepair/runner.py}, \texttt{make\_norm\_matcher}); it never copies the reference direction.

\section{Tuning}
\label{app:tuning}
Online training accuracy on the development seeds (1000, 1001) for each learning rate. All selections
were interior grid points, so no edge extension was used.
\begin{center}\scriptsize
\begin{tabular}{llcl}
\toprule
Condition & Family & selected LR & dev online acc.\ by LR \\
\midrule
RT & \texttt{adam} & 0.003 & 0.001: 0.451, 0.003: 0.544, 0.01: 0.519, 0.03: 0.400 \\
RT & \texttt{adam\_eqbeta} & 0.003 & 0.001: 0.414, 0.003: 0.478, 0.01: 0.452, 0.03: 0.384 \\
RT & \texttt{sgd} & 0.01 & 0.003: 0.492, 0.01: 0.513, 0.03: 0.451, 0.1: 0.320 \\
LP & \texttt{adam} & 0.003 & 0.001: 0.273, 0.003: 0.361, 0.01: 0.347, 0.03: 0.237 \\
LP & \texttt{adam\_eqbeta} & 0.003 & 0.001: 0.228, 0.003: 0.269, 0.01: 0.241, 0.03: 0.192 \\
LP & \texttt{sgd} & 0.01 & 0.003: 0.303, 0.01: 0.347, 0.03: 0.309, 0.1: 0.189 \\
\bottomrule
\end{tabular}

\end{center}

\section{Descriptive statistics of the intervention branches}
\label{app:descriptive}
Mean update norm over the first 10 probe steps ($\bar{\|u\|}_{1:10}$), relative parameter change after
200 probe steps, linear CKA of hidden features to the starting point on 64 fixed inputs, and accuracy
drop on task 49 (mean of three seeds). These are descriptive only.
\begin{center}\scriptsize\setlength{\tabcolsep}{3pt}
\begin{tabular}{l cccc cccc}
\toprule
 & \multicolumn{4}{c}{random-teacher ($C{=}4$)} & \multicolumn{4}{c}{label-permutation ($C{=}10$)} \\
\cmidrule(lr){2-5}\cmidrule(lr){6-9}
Arm & $\bar{\|u\|}_{1:10}$ & rel.\ $\Delta\theta$ & CKA & old drop & $\bar{\|u\|}_{1:10}$ & rel.\ $\Delta\theta$ & CKA & old drop \\
\midrule
\texttt{keep\_all} & 0.0311 & 0.38 & 0.703 & 0.441 & 0.0347 & 0.342 & 0.773 & 0.398 \\
\texttt{reset\_m} & 0.0726 & 0.391 & 0.689 & 0.449 & 0.0808 & 0.349 & 0.786 & 0.408 \\
\texttt{reset\_v} & 16.4 & 15.1 & 0.351 & 0.418 & 330 & 274 & 0.159 & 0.365 \\
\texttt{reset\_both} & 0.049 & 0.448 & 0.756 & 0.444 & 0.0403 & 0.386 & 0.845 & 0.400 \\
\texttt{reset\_t} & 0.00572 & 0.169 & 0.925 & 0.143 & 0.00632 & 0.16 & 0.903 & 0.109 \\
\texttt{reset\_v}, shared-keep$^\dagger$ & 17.3 & 15.7 & 0.294 & 0.487 & 331 & 273 & 0.160 & 0.368 \\
\texttt{reset\_both}, shared-keep$^\dagger$ & 0.249 & 0.859 & 0.471 & 0.538 & 0.209 & 0.711 & 0.723 & 0.426 \\
\bottomrule
\end{tabular}

\end{center}

\section{Planned experiments (NOT RUN)}
\label{app:planned}
None of the following has been run, and no result is implied:
\begin{enumerate}
\item a learning-rate-matched control for the shared-keep artifact arms, using the schedule of
Section~\ref{sec:analysis}(ii);
\item tuning the learning rate for the probe objective, to test whether the late-minus-early gap
persists;
\item more probe tasks per condition, fresh test seeds that have not been inspected, and a real-data
stream;
\item coordinate-level traces at the late checkpoint, to test the near-zero-gradient explanation of
\texttt{reset\_v};
\item any boundary-free trigger or predictor of intervention benefit (deliberately out of scope).
\end{enumerate}

\section{Correction history}
\label{app:corrections}
An earlier internal report (\texttt{docs/STAGE1\_REPORT.md}) stated that \texttt{reset\_t} equals a
per-tensor learning-rate rescaling ``exactly'' and that the first-step ratio is 0.3162. Both statements
are corrected here. The equivalence holds only for $\eps=0$ and coordinates with $v_i>0$, and 0.3162 is
the $r=1$, $t\to\infty$ approximation of $\kappa(r,t)$ in \eqref{eq:ratio}. The earlier report also
did not distinguish one-step equivalence from trajectory equivalence. The observed ``0'' is an exact tie
in accuracy while continuous quantities differ at about \TrajLossDiff{} (Appendix~\ref{app:trace}). The
original report is kept unchanged in the repository, with a pointer to this correction.
